% GENERATED FILE. Edit canonical lessons, then run npm run generate:books.
% canonical-source-hash: fnv1a64:bcf1cee7b65c6da3
% canonical-lessons: ES-C557-luto, ES-C557-esquela, ES-C557-pesame, ES-C557-ataud, ES-C557-herencia

\chapter{De luto}
\label{ch:es-de-luto}
\hlchaptermodality{557}

\begin{chapteropening}
\textbf{By the end of this chapter:} \emph{I can talk about mourning: a death notice, condolences, the coffin and what someone inherited.}

Complete the last lesson of chapter 557: I can talk about mourning: a death notice, condolences, the coffin and what someone inherited.
\end{chapteropening}

\section[el luto]{el luto — mourning}

\label{lesson:ES-C557-luto}



\begin{quote}
Before the new one: say the Spanish for the grave, then the Spanish for the cemetery.
\end{quote}

\subsection*{You'll want to know: el luto}
\textbf{el luto} — \textquotedblleft{}mourning\textquotedblright{}.

\begin{quote}
Toda la familia iba de luto.
\end{quote}

\begin{quote}
\emph{The whole family was in mourning.}
\end{quote}

\emph{Ir de luto} — to wear black for mourning.

\begin{cousinweb}
\emph{luto} is inherited from Latin \emph{lūctus} 'grief, mourning' (EtymDB), built on \emph{lūgēre} 'to mourn'.
\end{cousinweb}

\begin{grammarlens}[title={What you've built}]
One more word for telling, explaining and describing at B1.
\end{grammarlens}

\subsection*{Your turn}
\begin{itemize}
\raggedright
  \item \emph{Say it:} \emph{el luto}
  \item \emph{Say it:} \emph{el luto}, once more
  \item \emph{Say it:} a sentence of your own with \emph{el luto}
  \item \emph{Recall:} say the Spanish for the grave, then the Spanish for the cemetery, then say \emph{el luto} again
\end{itemize}

\begin{tcolorbox}[breakable,colback=teal!4,colframe=teal!35!black,title={Before you move on}]
What does el luto mean? (\textquotedblleft{}Mourning\textquotedblright{}.) Say it once more.
\end{tcolorbox}
\section[la esquela]{la esquela — the death notice}

\label{lesson:ES-C557-esquela}



\begin{quote}
Before the new one: say the Spanish for the cemetery, then the Spanish for mourning.
\end{quote}

\subsection*{You'll want to know: la esquela}
\textbf{la esquela} — \textquotedblleft{}the death notice\textquotedblright{}.

\begin{quote}
Leí su esquela en el periódico.
\end{quote}

\begin{quote}
\emph{I read his death notice in the paper.}
\end{quote}

A notice in the paper or on a wall telling people about a death and the funeral.

\begin{cousinweb}
\emph{esquela} is usually traced to Latin \emph{schedula} 'a small sheet of paper' (Lewis \& Short), first any short note and later the notice of a death.
\end{cousinweb}

\begin{grammarlens}[title={What you've built}]
One more word for telling, explaining and describing at B1.
\end{grammarlens}

\subsection*{Your turn}
\begin{itemize}
\raggedright
  \item \emph{Say it:} \emph{la esquela}
  \item \emph{Say it:} \emph{la esquela}, once more
  \item \emph{Say it:} a sentence of your own with \emph{la esquela}
  \item \emph{Recall:} say the Spanish for the cemetery, then the Spanish for mourning, then say \emph{la esquela} again
\end{itemize}

\begin{tcolorbox}[breakable,colback=teal!4,colframe=teal!35!black,title={Before you move on}]
What does la esquela mean? (\textquotedblleft{}The death notice\textquotedblright{}.) Say it once more.
\end{tcolorbox}
\section[el pésame]{el pésame — the condolences}

\label{lesson:ES-C557-pesame}



\begin{quote}
Before the new one: say the Spanish for mourning, then the Spanish for the death notice.
\end{quote}

\subsection*{You'll want to know: el pésame}
\textbf{el pésame} — \textquotedblleft{}the condolences\textquotedblright{}.

\begin{quote}
Le di el pésame a la viuda.
\end{quote}

\begin{quote}
\emph{I offered my condolences to the widow.}
\end{quote}

\emph{Dar el pésame} — to offer condolences. \emph{Mi más sentido pésame} — my deepest sympathy.

\begin{cousinweb}
\emph{pésame} is built on the phrase \emph{me pesa} 'it weighs on me, I am sorry', from \emph{pesar} 'to weigh', which comes from Latin \emph{pēnsāre} 'to weigh', built on \emph{pendere} 'to weigh, to hang' (Lewis \& Short).
\end{cousinweb}

\begin{grammarlens}[title={What you've built}]
One more word for telling, explaining and describing at B1.
\end{grammarlens}

\subsection*{Your turn}
\begin{itemize}
\raggedright
  \item \emph{Say it:} \emph{el pésame}
  \item \emph{Say it:} \emph{el pésame}, once more
  \item \emph{Say it:} a sentence of your own with \emph{el pésame}
  \item \emph{Recall:} say the Spanish for mourning, then the Spanish for the death notice, then say \emph{el pésame} again
\end{itemize}

\begin{tcolorbox}[breakable,colback=teal!4,colframe=teal!35!black,title={Before you move on}]
What does el pésame mean? (\textquotedblleft{}The condolences\textquotedblright{}.) Say it once more.
\end{tcolorbox}
\section[el ataúd]{el ataúd — the coffin}

\label{lesson:ES-C557-ataud}



\begin{quote}
Before the new one: say the Spanish for the death notice, then the Spanish for the condolences.
\end{quote}

\subsection*{You'll want to know: el ataúd}
\textbf{el ataúd} — \textquotedblleft{}the coffin\textquotedblright{}.

\begin{quote}
Seis amigos llevaron el ataúd.
\end{quote}

\begin{quote}
\emph{Six friends carried the coffin.}
\end{quote}

In Latin America \emph{el féretro} is also used.

\begin{cousinweb}
\emph{ataúd} was borrowed from Arabic \emph{tābūt} 'a box, a coffin' (EtymDB), with the Arabic article joined to the front.
\end{cousinweb}

\begin{grammarlens}[title={What you've built}]
One more word for telling, explaining and describing at B1.
\end{grammarlens}

\subsection*{Your turn}
\begin{itemize}
\raggedright
  \item \emph{Say it:} \emph{el ataúd}
  \item \emph{Say it:} \emph{el ataúd}, once more
  \item \emph{Say it:} a sentence of your own with \emph{el ataúd}
  \item \emph{Recall:} say the Spanish for the death notice, then the Spanish for the condolences, then say \emph{el ataúd} again
\end{itemize}

\begin{tcolorbox}[breakable,colback=teal!4,colframe=teal!35!black,title={Before you move on}]
What does el ataúd mean? (\textquotedblleft{}The coffin\textquotedblright{}.) Say it once more.
\end{tcolorbox}
\section[la herencia]{la herencia — the inheritance}

\label{lesson:ES-C557-herencia}



\begin{quote}
Before the new one: say the Spanish for the condolences, then the Spanish for the coffin.
\end{quote}

\subsection*{You'll want to know: la herencia}
\textbf{la herencia} — \textquotedblleft{}the inheritance\textquotedblright{}.

\begin{quote}
Recibió una casa en herencia.
\end{quote}

\begin{quote}
\emph{He inherited a house.}
\end{quote}

\emph{Heredar} is to inherit.

\begin{cousinweb}
\emph{herencia} is usually traced to Latin \emph{haerentia} 'things that cling, belongings', from \emph{haerēre} 'to stick'. It was later pulled towards \emph{heredar} 'to inherit', from Latin \emph{hērēs} 'heir'.
\end{cousinweb}

\begin{grammarlens}[title={What you've built}]
One more word for telling, explaining and describing at B1.
\end{grammarlens}

\subsection*{Your turn}
\begin{itemize}
\raggedright
  \item \emph{Say it:} \emph{la herencia}
  \item \emph{Say it:} \emph{la herencia}, once more
  \item \emph{Say it:} a sentence of your own with \emph{la herencia}
  \item \emph{Recall:} say the Spanish for the condolences, then the Spanish for the coffin, then say \emph{la herencia} again
\end{itemize}

\begin{tcolorbox}[breakable,colback=teal!4,colframe=teal!35!black,title={Before you move on}]
What does la herencia mean? (\textquotedblleft{}The inheritance\textquotedblright{}.) Say it once more.
\end{tcolorbox}
